\begin{center}
	\large{\textbf{EINSTEIN AND EINSTEIN-MAXWELL GRAVITY IN 2+1 DIMENSIONS AS A CHERN-SIMONS GAUGE THEORY}}
\end{center}

\vspace{1em}

\noindent
\begin{tabular}{@{}lcl@{}}
	Name       & : & Agung Sedayu Septiawan \\
	NRP        & : & 5001221075 \\
	Department & : & Physics \\
	Advisor    & : & Dr. rer. nat. Bintoro Anang Subagyo \\
\end{tabular}

\vspace{1em}

\noindent\textbf{Abstract}

\par This study formulates Einstein gravity in 2+1 dimensions as a Chern-Simons gauge theory with gauge group $\mathrm{SO}(2,2)$ and applies the formalism to BTZ black holes and the Einstein-Maxwell extension. Through the Witten ansatz, the spin connection and dreibein are unified into a single gauge connection valued in the Lie algebra $\mathfrak{so}(2,2)$, which is constructed from first principles; variation of the Chern-Simons action exactly reproduces the torsion-free condition and the Einstein field equations with negative cosmological constant, establishing the equivalence of the two formulations. The BTZ black hole solution and its event horizons are obtained analytically in the dreibein formalism, with the non-rotating limit analyzed as a special case that reduces the metric to diagonal form. The Einstein-Maxwell extension is carried out by adding the Maxwell action as a matter source; the Maxwell equations are solved in full generality and the electromagnetic energy-momentum tensor is computed explicitly in dreibein notation. A consistency analysis reveals that a geometric structural identity imposed by the BTZ ansatz clashes with the asymmetry of the electromagnetic energy-momentum tensor, so that a simultaneously charged and rotating solution does not exist in pure Einstein-Maxwell theory in 2+1 dimensions.

\vspace{1em}

\noindent\textbf{Keywords}: \textit{2+1 gravity, Chern-Simons, BTZ, Einstein-Maxwell}

\newpage
\thispagestyle{empty}
\vspace*{\fill}
    \begin{center}
	Halaman ini sengaja dikosongkan
    \end{center}
\vspace*{\fill}
\pagebreak
